\section{Bryan and Steinberg Composition at
\texorpdfstring{$N\in\{4,6\}$}{N in \{4,6\}}}%
\label{wn:sec:bryan-steinberg-N46}

\providecommand{\kBKM}{\kappa_{\mathrm{BKM}}}
\providecommand{\kch}{\kappa_{\mathrm{ch}}}
\providecommand{\kcat}{\kappa_{\mathrm{cat}}}
\providecommand{\kfib}{\kappa_{\mathrm{fiber}}}
\providecommand{\DTred}{Z^{\mathrm{red}}_{\mathrm{DT}}}
\providecommand{\FCHL}{F^{\mathrm{CHL}}}
\providecommand{\Z}{\mathbb{Z}}

\paragraph{Primitive CHL denominators.}
Let \(S\) be a K3 surface, let \(E\) be an elliptic curve, and let
\(g_N\) be a Mukai-symplectic automorphism of \(S\) of Mathieu class
\(4B\) or \(6A\). The primitive CHL denominator is governed by the
Borcherds weight formula
\[
\kBKM(\Phi_N)=\frac{c_N(0)}{2}.
\]
For the two orders considered here the Gritsenko and Nikulin
paramodular tower gives
\[
c_4(0)=4,\qquad \kBKM(\Phi_4)=2,
\]
and
\[
c_6(0)=2,\qquad \kBKM(\Phi_6)=1.
\]
The \(6A\) frame shape is
\[
1^2\,2^2\,3^2\,6^2,
\]
which has degree \(2+4+6+12=24\). This is the Cheng, Duncan, and
Harvey umbral \(6D_4\) row with umbral group \(3.\mathrm{Sym}_6\) of
order \(2160\). The Niemeier lattice \(A_5^4D_4\) is the separate
Coxeter-six Niemeier entry attached to a different umbral character.

\paragraph{Gritsenko and Clery diagonal-divisor catalogue.}
The Gritsenko and Clery diagonal-divisor catalogue is indexed by
paramodular row, divisor, and multiplier system. Its twice-weight data
are
\[
\begin{array}{c@{\qquad}cccccccc}
\mathrm{row} & 1 & 2 & 3 & 4 & 5 & 6 & 7 & 8\\
2\,\mathrm{wt} & 10 & 4 & 6 & 2 & 4 & 1 & 3 & 2 .
\end{array}
\]
A comparison with the primitive CHL ladder requires an explicit row map
and normalisation map. The constants \(c_4(0)=4\) and \(c_6(0)=2\)
above are the primitive CHL constants.

\paragraph{Local crepant-resolution theorem.}
Bryan and Steinberg prove a PT/DT crepant-resolution formula for a
crepant resolution of the coarse space of a Calabi-Yau orbifold
satisfying hard Lefschetz. The fibrewise orbifold
\[
\mathcal X_N=[S\times E/\langle(g_N,1)\rangle]
\]
has local cyclic \(A\)-type charts \([\mathbb C^2/\Z_r]\times E\).
The stabilisers are finite abelian subgroups of \(\mathrm{SL}_3\): the
surface action preserves \(\omega_S\), the elliptic factor contributes
\(dz_E\), and the product form \(\omega_S\wedge dz_E\) is invariant.
Bryan, Cadman, and Young compute the corresponding cyclic orbifold
vertices. Jun Li degeneration supplies the standard fibrewise gluing
formalism.

\paragraph{Order four.}
For \(g_4\) of frame shape \(1^4\,2^2\,4^4\), Nikulin fixed-point data
give
\[
\#\mathrm{Fix}(g_4)=4,\qquad \#\mathrm{Fix}(g_4^2)=8.
\]
Thus \(S/\langle g_4\rangle\) has four stabiliser-\(\Z_4\) orbits and
two stabiliser-\(\Z_2\) orbits:
\[
S/\langle g_4\rangle\quad\hbox{has Du Val type}\quad 4A_3+2A_1 .
\]
The Euler check is
\[
e(S/\Z_4)=\frac{24+4+8+4}{4}=10,\qquad
10+4\cdot3+2\cdot1=24.
\]
The local Bryan, Cadman, and Young factors are indexed by the \(A_3\)
and \(A_1\) cyclic charts. The primitive denominator in the reduced
numerical identity has weight
\[
\mathrm{wt}\,\FCHL_4=\frac{c_4(0)}{2}=2.
\]

\paragraph{Order six.}
For \(g_6\) of frame shape \(1^2\,2^2\,3^2\,6^2\), Nikulin fixed-point
data give
\[
\#\mathrm{Fix}(g_6)=2,\qquad
\#\mathrm{Fix}(g_6^2)=6,\qquad
\#\mathrm{Fix}(g_6^3)=8.
\]
The two \(g_6\)-fixed points give two \(A_5\) charts. The four
remaining points fixed by \(\langle g_6^2\rangle\cong\Z_3\) form two
\(\Z_6\)-orbits and give two \(A_2\) charts. The six remaining points
fixed by \(\langle g_6^3\rangle\cong\Z_2\) form two \(\Z_6\)-orbits
and give two \(A_1\) charts. Hence
\[
S/\langle g_6\rangle\quad\hbox{has Du Val type}\quad
2A_5+2A_2+2A_1 .
\]
The Euler check is
\[
e(S/\Z_6)=\frac{24+2+6+8+6+2}{6}=8,\qquad
8+2\cdot5+2\cdot2+2\cdot1=24.
\]
The local Bryan, Cadman, and Young factors are indexed by the \(A_5\),
\(A_2\), and \(A_1\) cyclic charts. The primitive denominator in the
reduced numerical identity has weight
\[
\mathrm{wt}\,\FCHL_6=\frac{c_6(0)}{2}=1.
\]

\paragraph{Reduced Donaldson-Thomas identity.}
The numerical identity in the primitive CHL normalisation is
\[
\DTred(\mathcal X_N)(q,p,y)
=-\,C_N(q,y)\,{\FCHL_N(q,p,y)}^{-2},
\qquad
C_N(q,y)=2\phi^{(g_N)}_{0,1}(q,y),
\]
with \(N=4\) or \(N=6\). The theorem-level ingredients are the
Borcherds weights above, the Nikulin stabiliser-orbit calculation, the
local cyclic orbifold vertices, and the Bryan and Steinberg
crepant-resolution formula. The global reduced identity is conditional
on the degeneration-gluing comparison between the local vertex product,
the reduced \(K3\times E\) cosection, and the primitive CHL denominator
\(\FCHL_N\).

\paragraph{Motivic and K-theoretic refinements.}
Maulik and Toda construct the motivic Donaldson-Thomas class with
Behrend cosection for K3-fibred threefolds. The cosection descends to
\(\mathcal X_N\) because \(g_N\) preserves \(\omega_S\). The finite
abelian equivariant Behrend function is compatible with the
Kontsevich and Soibelman motivic wall-crossing formalism, so the
Euler specialisation agrees with the numerical local theory whenever
the degeneration-gluing comparison holds.

The untwined refined \(K3\times E\) theory supplies the reduced
\(\chi_y\)/Jacobi model through Oberdieck's refined stable-pair
formula. Its \(N=4,6\) orbifold extension requires compatibility
between the reduced \(E\)-fibred multiplicative structure, the
\(g_N\)-linearised obstruction theory, and the \(M_{24}\) twining
sectors. The K-theoretic orbifold refinement remains conjectural at
this point.

\paragraph{Separated invariants.}
The ambient \(K3\times E\) split consists of four independent
constructions:
\[
\begin{aligned}
\kcat(K3\times E)
  &=\chi(\mathcal O_{K3})\,\chi(\mathcal O_E)=2\cdot0=0,\\
\kch^{\mathrm{Heis}}(K3\times E)
  &=3,\\
\kBKM(\mathfrak g_{\Delta_5})
  &=\frac{c_1(0)}2=5,\\
\kfib(K3\times E)
  &=\mathrm{rk}\,H^*(K3,\Z)=24 .
\end{aligned}
\]
Let
\[
X_N=(S\times E)/\langle(g_N,t_N)\rangle,\qquad
t_N\in E[N]\ \hbox{of exact order}\ N,
\]
be the smooth diagonal CHL quotient. For \(N=4\), the corresponding
data used here are
\[
\begin{aligned}
\kBKM(\Phi_4)&=2,\\
\kcat(X_4)&=0,\\
\kch(X_4)&=\chi(\mathcal O_{X_4})=0,\\
\kfib^{(4)}&=14 .
\end{aligned}
\]
For \(N=6\), the corresponding data used here are
\[
\begin{aligned}
\kBKM(\Phi_6)&=1,\\
\kcat(X_6)&=0,\\
\kch(X_6)&=\chi(\mathcal O_{X_6})=0,\\
\kfib^{(6)}&=14 .
\end{aligned}
\]
The Borcherds weight, the categorical Euler characteristic, the chiral
Hodge supertrace, and the fibre Mukai rank are separate invariants.
